\documentclass[11pt]{article}

\usepackage[a4paper,margin=29mm]{geometry}
\usepackage{fontspec}
\usepackage{unicode-math}
\setmainfont{TeX Gyre Termes}
\setsansfont{TeX Gyre Heros}
\setmonofont{DejaVu Sans Mono}
\setmathfont{TeX Gyre Termes Math}
\newfontfamily\zhfont[BoldFont={Droid Sans Fallback}]{Droid Sans Fallback}
\usepackage[final]{microtype}
\usepackage{mathtools,amsthm}
\usepackage{enumitem}
\usepackage{ragged2e}
\usepackage[numbers,sort&compress]{natbib}
\usepackage{xcolor}
\usepackage{booktabs}
\usepackage{url}
\usepackage[colorlinks=true,linkcolor=blue!45!black,citecolor=blue!45!black,
  urlcolor=blue!45!black]{hyperref}

\newtheorem{theorem}{Theorem}[section]
\newtheorem{proposition}[theorem]{Proposition}
\newtheorem{lemma}[theorem]{Lemma}
\newtheorem{corollary}[theorem]{Corollary}
\newtheorem{remark}[theorem]{Remark}
\theoremstyle{definition}
\newtheorem{definition}[theorem]{Definition}

\newcommand{\Spec}{\operatorname{Spec}}
\newcommand{\Ext}{\operatorname{Ext}}
\newcommand{\Hom}{\operatorname{Hom}}
\newcommand{\Wbig}{W}
\newcommand{\Wrat}{W_{\mathrm{rat}}}
\newcommand{\Zpre}{\underline{\mathbb Z}}
\newcommand{\Zsh}{\mathcal Z}
\newcommand{\Wsh}{\mathcal W}
\newcommand{\Ksh}{\mathcal K}
\newcommand{\sh}{\sharp}
\newcommand{\latin}[1]{{\rmfamily #1}}
\newcommand{\ff}{\mathrm{ff}}
\newcommand{\fppf}{\mathrm{fppf}}
\newcommand{\muD}{\mu_d}
\newcommand{\eps}{\varepsilon}
\newcommand{\angles}[1]{\langle #1\rangle}

\title{\textbf{A Descent Obstruction to Verschiebung Lifts\\
on fppf and Finite-Flat Sites}}
\author{Liang Wang\textsuperscript{1}\\
\small \textsuperscript{1}School of Artificial Intelligence and Automation,\\
\small Huazhong University of Science and Technology,\\
\small Luoyu Road 1037, 430070, Hubei, P.R. China\\
\small Contact: \texttt{wangliang.f@gmail.com}}
\date{Draft of 25 August 2026}

\begin{document}
\maketitle

\begin{abstract}
Let \(\omega\colon\mathcal Z\twoheadrightarrow\mathcal W\) be the
epimorphism from the sheafified reduced monoid algebra to the rational
big-Witt sheaf on a universe-small absolute site of noetherian affine
schemes.  We prove that, for every integer \(N>1\), the additive
Verschiebung \(V_N(f)(T)=f(T^N)\) admits no additive sheaf lift through
\(\omega\), for either the fppf topology or the finite-flat topology.
The obstruction is explicit.  After writing \(N=q^a d\) and passing from
\(k[x]\) to the finite-free root cover \(k[s]\), \(x=s^N\), injectivity
over the Dedekind domain \(k[s]\) forces a unique local preimage.  Its two
pullbacks disagree on the double overlap; specialization to
\(k[\varepsilon]/(\varepsilon^N)\) detects a nonzero kernel section.
Consequently, for the extension
\[
0\longrightarrow\mathcal K\longrightarrow\mathcal Z
 \xrightarrow{\omega}\mathcal W\longrightarrow0,
\]
one has \(u_*e\ne V_N^*e\) for every endomorphism
\(u\colon\mathcal K\to\mathcal K\).  The same example shows that the
sectionwise Dedekind-ring assertion in Corollary~4.6 of Deninger's
version-1 preprint, read literally, requires correction: sheaf
epimorphy gives local rather than objectwise surjectivity.
This answers the cited lifting question at every nontrivial index.
\end{abstract}

\begin{center}
\begin{minipage}{0.94\textwidth}
{\zhfont\small\RaggedRight
中文摘要。设 \(\omega\colon\mathcal Z\twoheadrightarrow\mathcal W\)
为诺特仿射概形绝对小站点上的层满态射，\\
其定义域来自约化幺半群代数的层化，值域为有理大
\latin{Witt} 层。\\
本文证明：对每个整数 \(N>1\)，加法 \latin{Verschiebung}
\(V_N(f)(T)=f(T^N)\) 在 \latin{fppf} 拓扑和有限平坦拓扑上，\\
都不能通过 \(\omega\) 提升为 \(\mathcal Z\) 的加法层自同态。\\
证明给出显式下降障碍。写 \(N=q^a d\)，在有限自由根覆盖
\(k[x]\to k[s]\)、\(x\mapsto s^N\) 上，\\
\latin{Dedekind} 整环 \(k[s]\) 处的单射性迫使局部原像唯一；
其在双重交叠上的两个拉回并不相等。\\
把交叠特化到 \(k[\varepsilon]/(\varepsilon^N)\)，
可检测到一个非零核截面。\\
由此，对扩张类 \(e\) 及任意
\(u\colon\mathcal K\to\mathcal K\)，均有 \(u_*e\ne V_N^*e\)。\\
同一例子也说明 \latin{Deninger} 预印本第一版推论 \latin{4.6}
的逐截面 \latin{Dedekind} 环断言按字面理解时需要修正：\\
层的满态射只保证局部满射，并不保证每个对象上的截面满射。
}
\end{minipage}
\end{center}

\noindent\textbf{Keywords.}
rational Witt vectors; Verschiebung; fppf descent; finite-flat topology;
extensions of sheaves; reduced monoid algebra.

\medskip
\noindent\textbf{MSC 2020 (provisional).}
13F35; 14F20; 18G15.

\section{Introduction and main results}

The rational big-Witt functor admits a concrete presentation in terms of
formal differences of elements of a multiplicative monoid.  Deninger
recently studied the sheaf-theoretic form of this presentation, constructed
lifts of the big-Witt Frobenius operations, and asked whether the
Verschiebung operations can likewise be lifted to the sheafified reduced
monoid algebra
\cite[p.~25]{Deninger2025Rational}.  The question is deceptively close to
an elementary factorization problem.  Locally, the expression
\(1-xT^N\) can be split after adjoining an \(N\)-th root of \(x\).
Globally, however, the chosen factors must satisfy descent on the double
overlap.  We show that the local factorization forced by injectivity never
satisfies this condition when \(N>1\).

We work on a universe-small version \(\mathscr C\) of the absolute category
of noetherian affine schemes used in
\cite[Secs.~3--4]{Deninger2025Rational}.  It is equipped first with the
fppf topology.  The chosen universe is understood to contain all affine
schemes and fiber products occurring in the proof.  This formulation avoids
treating an arbitrary skeleton as though it were automatically closed under
the required base changes.  Write
\[
 \Zsh=\Zpre(\mathcal O)^\sh,\qquad
 \Wsh=\Wrat(\mathcal O),\qquad
 \omega\colon\Zsh\twoheadrightarrow\Wsh .
\]
Here \(\Zpre(A)\) is the reduced integral monoid algebra of the
multiplicative monoid of \(A\), and the superscript \({}^\sh\) denotes
abelian sheafification.  Deninger's Theorem~3.4 states that the rational
Witt presheaf is already an fpqc sheaf on this noetherian-affine owner
\cite[Thm.~3.4, p.~19]{Deninger2025Rational}; Proposition~4.3 gives the
displayed epimorphism of sheaves
\cite[Prop.~4.3, p.~21]{Deninger2025Rational}.

For \(N\geq1\), the additive big-Witt Verschiebung is
\[
 V_N\colon\Wsh\longrightarrow\Wsh,\qquad
 V_N(f)(T)=f(T^N).
\]
The word \emph{additive} refers to the Witt addition, represented by
multiplication of power series.  Our principal result is the following.

\begin{theorem}[All-index fppf nonlift]\label{thm:fppf}
For every integer \(N>1\), there is no morphism of abelian sheaves
\(\widetilde V_N\colon\Zsh\to\Zsh\) satisfying
\[
 \omega\circ\widetilde V_N=V_N\circ\omega
\]
on \(\mathscr C_{\fppf}\).  For \(N=1\), the identity morphism is a lift.
\end{theorem}

The topology called finite-flat in
\cite[Sec.~4]{Deninger2025Rational} is distinct from the fppf topology.
Here a finite-flat covering family of an affine scheme \(U\) means a
jointly surjective family \(\{U_i\to U\}\) in which every morphism is finite
and flat; on affines, the corresponding ring maps are finite locally free
and their spectra jointly cover \(U\).  We use the associated subcanonical
topology, so representable presheaves, and in particular the structure
sheaf used below, satisfy descent.  We therefore state the finite-flat
conclusion separately rather than deriving it formally from
Theorem~\ref{thm:fppf}.

\begin{theorem}[Finite-flat nonlift]\label{thm:ff}
On the same universe-small noetherian-affine category equipped with the
finite-flat topology, no additive lift of \(V_N\) through \(\omega\) exists
for any \(N>1\).
\end{theorem}

Both theorems are proved by one finite-free root cover, but each use of
injectivity and each sheaf-theoretic detector is checked for the topology
in question.  Choose a prime \(q\mid N\), write \(N=q^a d\) with
\((q,d)=1\), and take a finite field \(k\) of characteristic \(q\)
containing all \(d\)-th roots of unity.  Over
\[
 A=k[x]\longrightarrow B=k[s],\qquad x\longmapsto s^N,
\]
the image of any hypothetical lift at the generator \(x\) is forced to
restrict to
\[
 c(s)=q^a\sum_{\zeta\in\mu_d(k)}(\zeta s)\in\Zpre(B).
\]
On \(B\otimes_A B\), the two pullbacks of \(c(s)\) differ.  The
specialization \(s_1\mapsto\varepsilon,\ s_2\mapsto0\) to
\(k[\varepsilon]/(\varepsilon^N)\) turns the difference into a nonzero
section.  Thus the obstruction is an explicit failed descent condition,
not a dimension count or a formal appeal to an uncomputed cohomology
group.

There is also a useful extension-theoretic formulation.  For
\(\tau\in\{\fppf,\ff\}\), regard \(\omega\colon\Zsh\to\Wsh\) as a morphism
in \(\mathrm{Ab}(\mathscr C_\tau)\), put
\(\Ksh_\tau=\ker(\omega)\), and let
\[
 e_\tau:\quad 0\longrightarrow\Ksh_\tau\longrightarrow\Zsh
 \xrightarrow{\omega}\Wsh\longrightarrow0
 \tag{1.1}\label{eq:extension}
\]
denote its class in
\(\Ext^1_{\mathrm{Ab}(\mathscr C_\tau)}(\Wsh,\Ksh_\tau)\).
When one topology is fixed below, the unadorned symbols \(\Ksh\) and \(e\)
abbreviate \(\Ksh_\tau\) and \(e_\tau\), respectively.  Standard
pushout--pullback functoriality then gives the following consequence.

\begin{corollary}[Extension obstruction]\label{cor:ext}
For \(\tau\in\{\fppf,\ff\}\), \(N>1\), and every endomorphism
\(u\colon\Ksh_\tau\to\Ksh_\tau\) in
\(\mathrm{Ab}(\mathscr C_\tau)\),
\[
 u_*e_\tau\ne V_N^*e_\tau
 \qquad\text{in }
 \Ext^1_{\mathrm{Ab}(\mathscr C_\tau)}(\Wsh,\Ksh_\tau).
\]
In particular, \(e_\tau\ne0\) and \(V_N^*e_\tau\ne0\) for both values of
\(\tau\).
\end{corollary}

The proof also isolates a distinction relevant to a statement in the
source.  Corollary~4.6 of
\cite[p.~23]{Deninger2025Rational}, in its version-1 formulation, asserts a
sectionwise identification over Dedekind rings for the finite-flat site.
For \(A=k[x]\), the rational Witt section \(1-xT^N\) in our construction is
locally in the image of \(\omega\) but has no global preimage.  Read in this
sectionwise sense, that assertion therefore requires correction.  The
issue is the passage from an epimorphism of sheaves to surjectivity on the
sections of a fixed object.  Proposition~4.3 remains the needed local
surjectivity statement, and Proposition~4.5 remains compatible with the
injectivity argument below.

The exact source is Deninger's version-1 preprint.
Proposition~4.3 gives the sheaf epimorphism and hence local preimages,
whereas Proposition~4.5 supplies the conditional injectivity used here
after the required integral refinements
\cite[Props.~4.3 and~4.5, pp.~21--23]{Deninger2025Rational}.
Corollary~4.7 gives a positive comparison only on certain finer,
non-subcanonical topologies
\cite[Cor.~4.7, p.~24]{Deninger2025Rational}.  The closest earlier
algebraic presentation result is Deninger--Mellit's explicit kernel
computation for a localized monoid-algebra map to truncated \(S\)-Witt
vectors \cite[Thm.~1.1]{DeningerMellit2019}.  Its quotient is different
and it does not treat sheafification or finite-flat/fppf descent.  The
Stacks references used below supply general local-exactness and
pushout--pullback formalism
\cite[Tags 03CN, 010I, and 06XP]{StacksProject}, not the arithmetic lift
or obstruction.

A bounded update completed on 25 August 2026 searched the current arXiv
record, version history, and API; DataCite; exact DOI/title queries in
OpenAlex; Crossref; and the relevant official publisher records.  Query
clusters combined the phrases \emph{rational Witt vectors},
\emph{Verschiebung}, \emph{fppf}, \emph{sheafification},
\emph{finite flat}, \emph{reduced monoid algebra}, and \emph{descent};
separate post-7-August-2025 queries screened broader
Witt--Verschiebung and Witt--descent hits.  Inclusion required either the
same sheafified reduced-monoid-algebra-to-rational-Witt owner and a
finite-flat/fppf lift or nonlift, or a proposition-level presentation or
descent comparison.  The exact-owner clusters returned only the source or
zero hits.  Deninger--Mellit was retained as the nearest presentation
result; the recent arithmetic-jet and p-typical associative-Witt hits, and
the broader spherical, KEnd/TR, Galois-lifting, and
Grothendieck--Witt hits, were excluded because they change the owner.
Within this declared scope, the source remained at version 1 and no direct
post-source solution to the finite-flat/fppf additive lifting question was
located.  This bounded negative result is not a claim of global priority.

Our contribution is correspondingly narrow: we answer the additive lifting
question for \(V_N\), compute a concrete all-index descent obstruction for
the stated sheaf epimorphism, and record its formal extension consequence.
We do not construct a lift of Frobenius, a compatible system of Witt
operations, or a ring endomorphism.

The proof isolates a reusable descent-obstruction template.  Let
\[
 e:\quad 0\longrightarrow K\longrightarrow Z
 \xrightarrow{p}W\longrightarrow0
\]
be an extension of abelian sheaves.  Fix an object \(U\), a source section
\(z_0\in Z(U)\), and an endomorphism \(v\colon W\to W\), and put
\(w=v(p(z_0))\in W(U)\).  Suppose that a cover \(f\colon V\to U\)
admits a unique section \(z_V\in Z(V)\) with \(p(z_V)=f^*w\), and that
the two pullbacks of \(z_V\) to \(V\times_UV\) differ by a nonzero
section of \(\ker(p)\).  Then \(w\) has no global preimage in \(Z(U)\):
any global preimage would restrict to \(z_V\) and hence would have equal
pullbacks.  For any \(u\colon K\to K\), a middle-object map inducing
\((u,v)\) would send \(z_0\) to a global preimage of
\(v(p(z_0))=w\).  Proposition~\ref{prop:extcriterion} therefore turns
this particular failed descent test into \(u_*e\ne v^*e\) for every
\(u\).  The relation between \(z_0\), \(v\), and \(w\) is essential:
failure of an unrelated target section to have a global preimage would not
by itself imply these Ext inequalities.

In the present argument, \(p=\omega\), \(z_0=(x)^\sh\),
\(v=V_N\), \(w=1-xT^N\), \(U=\Spec k[x]\), and
\(V=\Spec k[s]\).  The roots-of-unity product supplies the local
preimage, Dedekind injectivity makes it unique, finite freeness and
subcanonicity validate the cover and its restrictions on each site, and
the truncated-nilpotent big-Witt detector together with torsion-freeness
proves that the overlap difference is nonzero.  These are the arithmetic
and site-specific inputs; the preceding conditional implication is the
reusable categorical core.

The order of proof reflects these separations.  Section~\ref{sec:objects}
fixes the presheaf, sheaf, and section-level notation and identifies the
actual extension.  Section~\ref{sec:lemmas} proves the two detection
statements and the site-dependent Dedekind injectivity lemma.
Section~\ref{sec:obstruction} constructs the forced local factorization and
shows that it fails descent for all \(N>1\).
Section~\ref{sec:ext} gives the extension formulation, while
Section~\ref{sec:finiteflat} records the finite-flat consequence and the
source-sensitive correction.  The final section states the controls and
the limits of the result.

\section{Rational Witt sheaves and the extension}\label{sec:objects}

For a commutative unital ring \(A\), let \(\mathbb Z[A]\) be the integral
monoid algebra on the multiplicative monoid of \(A\).  We write
\[
 \Zpre(A)=\mathbb Z[A]/\mathbb Z(0).
\]
Thus an element is represented additively by a finite formal sum
\(\sum_i n_i(a_i)\), with the basis element \((0)\) set equal to zero.
This is the reduced monoid algebra in equation~(4), p.~3, of
\cite{Deninger2025Rational}.  The power-series formula for Verschiebung used
below, \(V_N(f)(T)=f(T^N)\), is equation~(20), p.~14, of the same source.
The map to rational big-Witt vectors is
\[
 \omega_A\!\left(\sum_i n_i(a_i)\right)
   =\prod_i(1-a_iT)^{n_i}.
 \tag{2.1}\label{eq:omega}
\]
The target \(\Wrat(A)\) is viewed as an additive group under multiplication
of its representing rational power series.  Formula
\eqref{eq:omega} is natural in \(A\), and the Teichm\"uller element
\([a]\) is represented by \(1-aT\).

Three levels of notation occur and should not be conflated.  The first is
the presheaf group \(\Zpre(A)\), in which a displayed finite sum is written.
The second is the sheaf section group
\(\Zsh(\Spec A)\), to which that finite sum has a canonical image.  The
third is a section represented only after choosing a cover, which need not
come from any element of \(\Zpre(A)\).  The main local candidate happens to
be defined already at the first level, whereas the hypothetical global
preimage is allowed to be an arbitrary section at the second level.  Thus
the contradiction does not depend on an unjustified identification
\(\Zpre(A)=\Zsh(\Spec A)\).
Whenever a formal sum is displayed below, its canonical image in the
corresponding sheaf section group is understood.

Let \(\Zpre(\mathcal O)\) and \(\Wrat(\mathcal O)\) be the corresponding
presheaves on \(\mathscr C\).  We retain the notation
\[
 \Zsh=\Zpre(\mathcal O)^\sh,\qquad
 \Wsh=\Wrat(\mathcal O).
\]
The second equality uses the already-sheaf result quoted above and is made
only on the noetherian-affine site under discussion.  We do not silently
extend that assertion to a larger affine site without sheafifying the
target there.

It is essential that \(\omega\) is an epimorphism \emph{in the category of
sheaves}.  For sheaves of abelian groups, this means that a target section
has preimages after passage to a cover.  It does not mean that
\[
 \Zsh(U)\longrightarrow\Wsh(U)
\]
is surjective for every object \(U\).  Equivalently, the cokernel sheaf
vanishes even though an objectwise cokernel presheaf can have nonzero
sections before sheafification; see the local exactness criterion in
\cite[Tag 03CN]{StacksProject}.  The distinction is precisely what permits
a local factorization of \(1-xT^N\) while obstructing its descent.

More concretely, if \(w\in\Wsh(\Spec A)\), epimorphy yields a covering
family \(A\to A_i\) and sections \(z_i\in\Zsh(\Spec A_i)\) with
\(\omega(z_i)=w|_{A_i}\).  It supplies neither a natural choice of the
\(z_i\) nor their equality over \(A_i\otimes_A A_j\).  By contrast, an
endomorphism \(\widetilde V_N\) applied to the global section
\((x)^\sh\) would give a global preimage before any cover was selected.
Its restrictions would agree on every overlap by functoriality.  The proof
therefore tests exactly the extra compatibility that a sheaf epimorphism
does not provide.

Objectwise define
\[
 K_0(A)=\ker\bigl(\Zpre(A)\longrightarrow\Wrat(A)\bigr).
\]
In concrete terms,
\[
 K_0(A)=
 \left\{\sum_i n_i(a_i):
   \prod_i(1-a_iT)^{n_i}=1\right\}.
 \tag{2.2}\label{eq:kernel}
\]
Abelian sheafification is exact, so
\[
 \Ksh=K_0^\sh=\ker(\omega).
\]
This yields the genuine short exact sequence \eqref{eq:extension} in the
abelian category \(\mathrm{Ab}(\mathscr C_\tau)\), for
\(\tau=\fppf\) or \(\tau=\ff\).  Nothing here asserts that
\(K_0(A)\) vanishes, or that taking sections preserves the right exactness
of \eqref{eq:extension}.

The operation \(V_N\) is natural because substitution \(T\mapsto T^N\)
commutes with change of coefficients.  On Teichm\"uller representatives,
\[
 V_N([a])=1-aT^N.
 \tag{2.3}\label{eq:versch}
\]
A lift in this paper always means a morphism of the underlying abelian
sheaves satisfying the commuting square.  Such a lift would preserve
\(\Ksh\): if \(z\in\Ksh\), then
\(\omega(\widetilde V_Nz)=V_N(\omega z)=0\).  Hence it would induce some
endomorphism \(u\) of the kernel.  No multiplicative compatibility is
assumed.

\section{Three descent lemmas}\label{sec:lemmas}

We collect the ingredients that make the overlap calculation survive
sheafification.  The first rules out the possibility that the integer
coefficient \(q^a\) kills the obstruction.

\begin{lemma}[Torsion-freeness]\label{lem:torsionfree}
For every \(m\geq1\), multiplication by \(m\) is a monomorphism of
\(\Zsh\), for either topology.
\end{lemma}

\begin{proof}
The group \(\mathbb Z[A]\) is free abelian on the underlying multiplicative
monoid of \(A\).  The summand generated by \((0)\) is free and split as an
abelian group, so \(\Zpre(A)=\mathbb Z[A]/\mathbb Z(0)\) is free abelian on
the nonzero elements of \(A\).  Multiplication by \(m\) is therefore
objectwise injective.  Exactness of sheafification of abelian presheaves
preserves this monomorphism.
\end{proof}

The sheaf conclusion is stronger than torsion-freeness at one selected
ring.  If a section has survived sheafification, no nonzero integer
multiple of it can disappear after a further covering refinement.  In the
main calculation the detector first proves \(y^\sh\ne0\), and this lemma
then proves \(q^a y^\sh\ne0\).  Reversing those two steps would fail,
because the big-Witt image of the multiple is deliberately the identity.

We next use the full big-Witt sheaf only as a detector.  Under its standard
power-series description, the Teichm\"uller map is injective because the
coefficient of \(T\) in \(1-aT\) recovers \(-a\).

\begin{lemma}[Big-Witt detector]\label{lem:detector}
Let \(\tau\) be either the fppf or finite-flat topology.  If a section
\(z\in\Zsh(U)\) has nonidentity image in the big-Witt sheaf
\(\Wbig(\mathcal O)(U)\), then \(z\ne0\).
\end{lemma}

\begin{proof}
Both topologies are subcanonical, and the big-Witt construction in its
power-series model is a sheaf on them.  The natural map
\(\Zsh\to\Wbig(\mathcal O)\) sends a formal sum to the product in
\eqref{eq:omega}.  A zero section must map to the identity power series.
The contrapositive gives the claim.  This is the same detection principle
used for the class over
\(\mathbb F_2[\varepsilon]/(\varepsilon^2)\) in Deninger's
Example~4.4 \cite[p.~22]{Deninger2025Rational}.
\end{proof}

Subcanonicity enters only in this detection step and in interpreting the
displayed restrictions without changing the structure sheaf.  The
rational Witt map cannot detect the final obstruction, since that
obstruction lies in its kernel by construction.  Passing temporarily to
the full big-Witt sheaf retains the lower-degree coefficient
\(\varepsilon^dT^d\), and torsion-freeness then transports nonvanishing
back to the rational-kernel section that is needed.

The final lemma forces uniqueness of the local preimage on the root cover.

\begin{lemma}[Dedekind injectivity]\label{lem:dedekind}
Let \(B\) be a Dedekind domain.  The map
\[
 \omega_B\colon\Zsh(\Spec B)\longrightarrow\Wsh(\Spec B)
\]
is injective on \(\mathscr C_{\fppf}\).  It is also injective when
\(\mathscr C\) is equipped with the finite-flat topology.
\end{lemma}

\begin{proof}
Deninger's Proposition~4.5 gives injectivity when the relevant covers admit
refinements by integral schemes
\cite[Prop.~4.5, pp.~22--23]{Deninger2025Rational}.  We verify that
hypothesis in the two cases used here.

First take an fppf covering of \(\Spec B\).  Because the target is
quasi-compact, retain a finite subcover, written on rings as finitely
presented flat \(B\)-algebras \(C_i\).  Each \(C_i\) is noetherian.
Replace it by its finitely many quotients \(C_i/\mathfrak p\), where
\(\mathfrak p\) runs through the minimal primes.  Flat going down implies
\(\mathfrak p\cap B=(0)\); compare
\cite[Tag 00HS]{StacksProject}.  Thus every \(C_i/\mathfrak p\) is a
domain and is torsion-free as a \(B\)-module.  Over a Dedekind domain,
torsion-free modules are flat
\cite[Tag 0AUW]{StacksProject}.  The quotient remains of finite
presentation over \(B\), and the irreducible components jointly cover
\(\Spec C_i\).  These domains therefore form an fppf refinement of the
original cover.

For clarity, the finite-presentation assertion uses no hidden excellence
hypothesis: since \(C_i\) is noetherian, each minimal prime is finitely
generated, so \(C_i/\mathfrak p\) is again a finitely presented
\(B\)-algebra.  The union of the closed subsets
\(\Spec(C_i/\mathfrak p)\) over all minimal primes is all of
\(\Spec C_i\).  After composing with the original covering maps, the
resulting family is therefore still jointly surjective over \(\Spec B\).
Zero components, if present, may simply be omitted.
All rings retained in this finite construction lie in the fixed
universe-small site chosen at the outset.

If the original cover is finite flat, each \(C_i\) is finite over \(B\).
The same minimal-prime quotients are finite and torsion-free over \(B\),
hence finite locally free.  They give a finite-flat domain refinement.
Proposition~4.5 now applies separately in both topologies.
\end{proof}

\begin{remark}\label{rem:scope-site}
Lemma~\ref{lem:dedekind} is needed only at the polynomial ring
\(B=k[s]\), which is a principal ideal domain.  The proof does not assert
that \(\omega\) is injective at arbitrary nonreduced objects; indeed, the
nonreduced specialization in Section~\ref{sec:obstruction} is where the
kernel is detected.
\end{remark}

\section{The all-index descent obstruction}\label{sec:obstruction}

We now give the common explicit calculation.  Its output is slightly
stronger than nonliftability: for every nontrivial index it produces a
rational Witt section over a Dedekind domain that is locally, but not
globally, in the image of \(\omega\).

\begin{proposition}[Failure of objectwise surjectivity]\label{prop:failure}
Fix \(N>1\), and let \(\tau\) be either the fppf topology or the
finite-flat topology.  There are a finite field \(k\), the Dedekind domain
\(A=k[x]\), and a section
\[
 w_N=1-xT^N\in\Wrat(A)
\]
such that
\[
 w_N\notin
 \operatorname{im}\bigl(\omega_A\colon
 \Zsh(\Spec A)\longrightarrow\Wsh(\Spec A)\bigr).
\]
The section \(w_N\) acquires a preimage on a finite-free faithfully flat
cover of rank \(N\).
\end{proposition}

\begin{proof}
Choose a prime divisor \(q\) of \(N\), and write
\[
 N=q^a d,\qquad a\geq1,\qquad (q,d)=1.
 \tag{4.1}\label{eq:decompose}
\]
There is a finite extension \(k/\mathbb F_q\) containing all \(d\)-th
roots of unity.  Since \(q\nmid d\), the group
\(\mu_d(k)\) has exactly \(d\) elements.  Set
\[
 A=k[x],\qquad B=k[s],\qquad A\longrightarrow B,\quad x\longmapsto s^N.
 \tag{4.2}\label{eq:rootcover}
\]
As an \(A\)-module, \(B\) is free with basis
\(1,s,\ldots,s^{N-1}\); hence \(\Spec B\to\Spec A\) is a finite-free
faithfully flat cover.  In particular, it is a cover for both topologies.

Define the formal section
\[
 c_N(s)=q^a\sum_{\zeta\in\mu_d(k)}(\zeta s)
       \in\Zpre(B)\longrightarrow\Zsh(\Spec B).
 \tag{4.3}\label{eq:candidate}
\]
The elementary identity
\[
 \prod_{\zeta\in\mu_d(k)}(1-\zeta u)=1-u^d
 \tag{4.4}\label{eq:roots}
\]
and the characteristic-\(q\) Frobenius identity give
\begin{align*}
 \omega_B(c_N(s))
  &=\prod_{\zeta\in\mu_d(k)}(1-\zeta sT)^{q^a}\\
  &=(1-s^dT^d)^{q^a}\\
  &=1-s^{dq^a}T^{dq^a}
   =1-s^NT^N.
 \tag{4.5}\label{eq:localimage}
\end{align*}
Thus \(c_N(s)\) is a local preimage of the pullback of \(w_N\).
Moreover it is the only local preimage in \(\Zsh(\Spec B)\), because
\(B=k[s]\) is a Dedekind domain and Lemma~\ref{lem:dedekind} makes
\(\omega_B\) injective.

Suppose, for contradiction, that a global section
\(z\in\Zsh(\Spec A)\) satisfies \(\omega_A(z)=w_N\).  Its restriction to
\(\Spec B\) has the same Witt image as \(c_N(s)\), so uniqueness forces
\[
 z|_B=c_N(s).
 \tag{4.6}\label{eq:forced}
\]
The double overlap is
\[
 R=B\otimes_A B
   =k[s_1,s_2]/(s_1^N-s_2^N).
 \tag{4.7}\label{eq:overlap}
\]
Because the two restrictions of a global section agree, equation
\eqref{eq:forced} would imply
\[
 c_N(s_1)=c_N(s_2)\qquad\text{in }\Zsh(\Spec R).
 \tag{4.8}\label{eq:descentneeded}
\]

To test this equality, put
\[
 D=k[\varepsilon]/(\varepsilon^N)
 \tag{4.9}\label{eq:D}
\]
and use the homomorphism \(R\to D\) defined by
\(s_1\mapsto\varepsilon\) and \(s_2\mapsto0\).  It is well defined since
\(\varepsilon^N=0\).  Under the induced restriction map, the difference
between the two sides of \eqref{eq:descentneeded} becomes
\[
 q^a y^\sh,\qquad
 y=\sum_{\zeta\in\mu_d(k)}(\zeta\varepsilon)\in\Zpre(D).
 \tag{4.10}\label{eq:y}
\]
We first detect \(y^\sh\), rather than its \(q^a\)-multiple, in the
big-Witt sheaf:
\[
 \prod_{\zeta\in\mu_d(k)}(1-\zeta\varepsilon T)
   =1-\varepsilon^dT^d\ne1.
 \tag{4.11}\label{eq:detecty}
\]
Indeed \(d<N\), because \(a\geq1\), so
\(\varepsilon^d\ne0\) in \(D\).  Lemma~\ref{lem:detector} therefore gives
\(y^\sh\ne0\).  Lemma~\ref{lem:torsionfree} then gives
\[
 q^a y^\sh\ne0.
 \tag{4.12}\label{eq:qnonzero}
\]
This contradicts \eqref{eq:descentneeded}.

For completeness, the nonzero section in \eqref{eq:qnonzero} lies in the
kernel of the rational Witt map:
\[
 \omega_D(q^a y^\sh)
  =(1-\varepsilon^dT^d)^{q^a}
  =1-\varepsilon^NT^N=1.
 \tag{4.13}\label{eq:inkernel}
\]
Thus the calculation identifies an actual nonzero kernel section.  Notice
the order of the argument: the big-Witt image detects
\(y^\sh\), while torsion-freeness detects its \(q^a\)-multiple; the
big-Witt image of that multiple is itself the identity.
\end{proof}

\subsection{How the arithmetic factors enter}\label{subsec:factors}

Each part of the decomposition \(N=q^a d\) has a distinct role.  The field
characteristic turns the coefficient \(q^a\) in the monoid algebra into a
\(q^a\)-th power on the Witt side.  The equality
\((1-u)^{q^a}=1-u^{q^a}\) then changes degree \(d\) into degree \(N\).
The prime-to-\(q\) factor is handled by the separable polynomial
\(X^d-1\): after a finite extension of \(\mathbb F_q\), all of its roots
are available and equation \eqref{eq:roots} packages them without
multiplicity.  No algebraic closure is needed, and the resulting field
remains finite.

The strict inequality \(d<N\) is equally important.  It ensures that
\(\varepsilon^d\) survives in
\(D=k[\varepsilon]/(\varepsilon^N)\), so the inner section is visible to
the big-Witt detector.  At the same time, multiplication by \(q^a\) raises
the detected power series to
\(1-\varepsilon^NT^N=1\), placing the final nonzero section in the rational
kernel.  Thus the same factorization simultaneously creates visibility
before multiplication and vanishing after multiplication.

Finally, the cover in \eqref{eq:rootcover} is not chosen merely for
convenience.  Its finite freeness places it in both sites, while its source
\(k[s]\) is a PID, where Lemma~\ref{lem:dedekind} forces uniqueness.
These two properties let a single overlap calculation address the two
topologies while keeping their injectivity proofs logically separate.

The construction is uniform in the index but intentionally not in the
coefficient field.  For each \(N\), a prime divisor \(q\) and a finite
extension of \(\mathbb F_q\) containing \(\mu_d\) are selected.  This is
enough to rule out a natural transformation on the absolute site: such a
transformation must act on every object, so one counterexample object for
each index is decisive.  No single coefficient ring is claimed to detect
all indices simultaneously.

Two possible escape routes are also excluded by the form of the proof.
First, the section \(z\), if it existed, could be an arbitrary sheaf
section; it is not assumed to be represented by one formal sum over
\(A\).  Only after restricting to \(B\) does injectivity identify it with
the displayed presheaf section \(c_N(s)\).  Second, a different covering
family cannot avoid the contradiction.  Any global section produced using
such a family would still restrict to the fixed root cover and would again
be forced to equal \(c_N(s)\), so it would have the same nonzero overlap
difference.

The direction of specialization is similarly important.  The ring map
\(R\to D\) defines a morphism \(\Spec D\to\Spec R\), hence a restriction of
sheaf sections from the overlap to the truncated ring.  Nonvanishing after
this restriction proves nonvanishing before it.  No claim is made that
\(D\) covers \(R\), and none is needed: the specialization is a detector
of inequality, not a descent cover.
It is well defined precisely because both \(s_1^N\) and \(s_2^N\) map to
zero.  Every ring in this detection chain remains noetherian and affine.

The obstruction can be retained directly on the double overlap.  Define
\[
 \delta_N=c_N(s_1)-c_N(s_2)\in\Zsh(\Spec R).
 \tag{4.14}\label{eq:delta}
\]
Since \(s_1^N=s_2^N\) in \(R\), equation \eqref{eq:localimage} shows that
\(\delta_N\in\Ksh(\Spec R)\).  Its image under \(R\to D\) is the nonzero
section \(q^a y^\sh\).  Hence \(\delta_N\ne0\).  This gives a compact
description of the failure: the only possible local preimage has a
nonvanishing difference on the first overlap.

\begin{proof}[Proof of Theorems~\ref{thm:fppf} and \ref{thm:ff}]
Fix one of the two topologies and suppose that
\(\widetilde V_N\) is a lift for some \(N>1\).  With \(A=k[x]\) as in
Proposition~\ref{prop:failure}, apply the lift to the section
\((x)^\sh\in\Zsh(\Spec A)\).  Naturality and the commuting-square
condition give
\[
 \omega_A\bigl(\widetilde V_N((x)^\sh)\bigr)
  =V_N(1-xT)=1-xT^N=w_N.
\]
This is a global preimage excluded by Proposition~\ref{prop:failure}.
The argument proves the fppf theorem.  Separately, the root morphism
\eqref{eq:rootcover} is finite free, the finite-flat form of
Lemma~\ref{lem:dedekind} supplies uniqueness, and the finite-flat topology
is subcanonical; the same calculation therefore proves the finite-flat
theorem without transferring a conclusion between sites.  If \(N=1\),
\(\widetilde V_1=\operatorname{id}_{\Zsh}\) satisfies the required
identity.
\end{proof}

\subsection{The first index as a source control}

When \(N=2\), take \(q=2\), \(a=1\), \(d=1\), and \(k=\mathbb F_2\).
Then
\[
 A=\mathbb F_2[x],\qquad B=\mathbb F_2[s],\qquad x=s^2,
\]
and the unique local preimage is \(c_2(s)=2(s)\).  The overlap is
\[
 \mathbb F_2[s_1,s_2]/(s_1^2-s_2^2)
  =\mathbb F_2[s_1,s_2]/((s_1-s_2)^2).
\]
Specialization to
\(\mathbb F_2[\varepsilon]/(\varepsilon^2)\) gives the nonzero class
\(2(\varepsilon)^\sh\) appearing in
Deninger's Example~4.4 \cite[p.~22]{Deninger2025Rational}.  The general
proof may therefore be viewed as an all-index root-of-unity expansion of
that first obstruction.

\section{The extension-theoretic formulation}\label{sec:ext}

For each \(\tau\in\{\fppf,\ff\}\), we relate the concrete failure to the
topology-indexed extension class \(e_\tau\) in
\eqref{eq:extension}.  The required statement is formal in the abelian
category \(\mathrm{Ab}(\mathscr C_\tau)\), but making the two functorial
directions explicit prevents a common ambiguity.

\begin{proposition}[Pushout--pullback criterion]\label{prop:extcriterion}
Let
\[
 e:\quad0\longrightarrow K\longrightarrow Z
   \xrightarrow{p}W\longrightarrow0
\]
be an extension in an abelian category.  Given
\(u\colon K\to K\) and \(v\colon W\to W\), there exists a morphism
\(\widetilde v\colon Z\to Z\) inducing \(u\) on \(K\), inducing \(v\) on
\(W\), and satisfying \(p\widetilde v=vp\), if and only if
\[
 u_*e=v^*e\qquad\text{in }\Ext^1(W,K).
 \tag{5.1}\label{eq:criterion}
\]
When nonempty, the set of such middle-object morphisms is a torsor under
\(\Hom(W,K)\).
\end{proposition}

\begin{proof}
Push out the first copy of \(e\) along \(u\) and pull back the second copy
along \(v\).  A morphism of the original extensions inducing the prescribed
end maps produces an equivalence between these two extensions.  Conversely,
an equivalence between the pushout and pullback extensions, composed with
their universal morphisms, produces a middle-object map with end maps
\((u,v)\).  This is the usual functoriality of extensions
\cite[Tag 010I]{StacksProject}; identifying Yoneda extensions with
\(\Ext^1\) gives \eqref{eq:criterion}
\cite[Tag 06XP]{StacksProject}.

If \(\widetilde v_1\) and \(\widetilde v_2\) have the same end maps, their
difference is zero on \(K\) and has image in \(K\).  It therefore factors
uniquely as
\[
 Z\xrightarrow{p}W\xrightarrow{h}K\longrightarrow Z
\]
for some \(h\in\Hom(W,K)\).  Adding such a factorization acts freely and
transitively on the set of middle-object maps.
\end{proof}

\begin{proof}[Proof of Corollary~\ref{cor:ext}]
Fix \(\tau\in\{\fppf,\ff\}\).  Apply
Proposition~\ref{prop:extcriterion} in
\(\mathrm{Ab}(\mathscr C_\tau)\) to the extension \(e_\tau\) in
\eqref{eq:extension}, with \(v=V_N\).  Equality
\(u_*e_\tau=V_N^*e_\tau\) would produce an additive middle-object
morphism \(\Zsh\to\Zsh\) on \(\mathscr C_\tau\) lifting \(V_N\), contrary
to Theorem~\ref{thm:fppf} when \(\tau=\fppf\) and to
Theorem~\ref{thm:ff} when \(\tau=\ff\).  Hence the two classes are unequal
for every \(u\colon\Ksh_\tau\to\Ksh_\tau\).  In the final paragraph of
this proof, the unadorned \(e\) and \(\Ksh\) abbreviate \(e_\tau\) and
\(\Ksh_\tau\).

Taking \(u=0\), the pushout \(0_*e\) is the zero extension class.  The
inequality therefore shows \(V_N^*e\ne0\).  Also \(e\ne0\): if
\eqref{eq:extension} split, the splitting would make \(\omega\) surjective
on the sections of every object, contradicting
Proposition~\ref{prop:failure}.  Equivalently, a chosen decomposition
\(Z\simeq K\oplus W\) would allow the middle map \(0\oplus V_N\).
\end{proof}

\begin{remark}[What the overlap calculation does not claim]
The section \(\delta_N\) in \eqref{eq:delta} is a necessary-descent
detector for a specific finite-free cover.  We do not claim that the
associated \v Cech complex computes the full sheaf
\(\Ext^1(\Wsh,\Ksh)\).  The passage from the explicit nonlift to the Ext
inequality uses Proposition~\ref{prop:extcriterion}, not a comparison
theorem between \v Cech and derived cohomology.
\end{remark}

\subsection{The connecting-section viewpoint}

There is a second way to state the nonvanishing that remains below the
level of a full \v Cech-to-derived comparison.  Apply the section functor
over \(\Spec A\) to the short exact sequence
\eqref{eq:extension}.  Its long exact cohomology sequence contains a
connecting morphism
\[
 \partial_A\colon
 \Wsh(\Spec A)\longrightarrow
 H^1_\tau(\Spec A,\Ksh).
 \tag{5.2}\label{eq:connecting}
\]
Exactness says that \(\partial_A(w_N)=0\) precisely when \(w_N\) has a
global preimage in \(\Zsh(\Spec A)\).  Proposition~\ref{prop:failure}
therefore gives
\[
 \partial_A(1-xT^N)\ne0.
 \tag{5.3}\label{eq:connectingnonzero}
\]
The section \(\delta_N\) is an explicit witness to the failed first
descent condition for the chosen root cover.  Equation
\eqref{eq:connectingnonzero}, however, follows from the long exact
sequence and objectwise nonsurjectivity, so it does not require the
additional assertion that this particular \v Cech cocycle computes all of
\(H^1_\tau\), much less all of \(\Ext^1\).

\section{The finite-flat site and the Dedekind-section assertion}
\label{sec:finiteflat}

Although the same equations occur in both topologies, the finite-flat
conclusion has its own site-dependent inputs.  The cover
\(k[x]\to k[s]\) is finite free; its double overlap remains an object of
the noetherian-affine site; the specialization to
\(k[\varepsilon]/(\varepsilon^N)\) is a valid restriction; and both the
rational and full big-Witt targets are sheaves for the topology in use.
Most importantly, the finite-flat half of
Lemma~\ref{lem:dedekind} refines a finite-flat cover of a Dedekind domain
by finite-flat integral domains.  These checks, rather than a change-of-site
implication, justify Theorem~\ref{thm:ff}.

We now spell out the consequence for the source statement.  On the
finite-flat site, take already the first index
\[
 A=\mathbb F_2[x],\qquad w=1-xT^2\in\Wrat(A).
 \tag{6.1}\label{eq:cor46example}
\]
The section \(w\) has the local preimage \(2(s)^\sh\) on the finite-free
cover \(x=s^2\), as Proposition~4.3 predicts.  Proposition~4.5, together
with the Dedekind-domain refinement in Lemma~\ref{lem:dedekind}, makes that
preimage unique over \(B=\mathbb F_2[s]\).  Its two restrictions disagree
on the overlap, as detected by \(2(\varepsilon)^\sh\ne0\).  Consequently
\[
 w\notin
 \operatorname{im}\bigl(
 \Zsh(\Spec\mathbb F_2[x])\longrightarrow
 \Wrat(\mathbb F_2[x])\bigr).
 \tag{6.2}\label{eq:notimage}
\]

Corollary~4.6 of the version-1 source states a sectionwise equality over
Dedekind rings for the finite-flat topology
\cite[p.~23]{Deninger2025Rational}.  Equation \eqref{eq:notimage} shows
that this equality does not hold as stated.  The proof in the source
appears to use the epimorphism of Proposition~4.3 as though it were
surjective on the section group at \(\Spec A\).  In a sheaf category, an
epimorphism supplies local preimages; it does not ensure that a chosen
local preimage descends.  Proposition~4.5 supplies the relevant
injectivity after the domain-refinement argument, but injectivity only
makes the failed local choice unique.  It cannot turn it into a global
section.

This correction is deliberately limited.  It does not affect the
epimorphism of sheaves in Proposition~4.3, the conditional injectivity
mechanism of Proposition~4.5, or the explicit noninjectivity phenomenon in
Example~4.4.  It concerns the objectwise-surjectivity step in the printed
Corollary~4.6.  The calculation is also independent of any response or
future revision by the source author.

\subsection{Robustness of the section counterexample}

The contradiction is not an artifact of having chosen an inconvenient
local preimage.  Suppose a global preimage of \(w_N\) existed, perhaps
constructed using a completely different covering family.  Restricting
that global section to the fixed root cover would still produce a preimage
of \(1-s^NT^N\) over \(B=k[s]\).  Dedekind injectivity would identify it
with \(c_N(s)\), whose overlap restrictions have already been proved
unequal.  Thus no alternative cover, refinement, or local factorization can
repair the global section.  The root cover is a test forced on every
putative global preimage, not an assumption about how such a preimage must
first be found.

Nor does the argument rely on \(w_N\) lying only in a larger completion.
It is the Verschiebung \(V_N([x])\) of the rational Teichm\"uller element
\([x]=1-xT\), so \(w_N=1-xT^N\) is a rational Witt vector by naturality of
\(V_N\).  All displayed products in the proof are finite.  The
specialization ring \(D\) is used only to test equality of sheaf sections;
it is not asserted to be Dedekind, and no injectivity claim is made there.
This separation is what allows nilpotents to expose the kernel without
weakening the uniqueness argument over \(B\).

The source inputs consequently have distinct logical roles.
Theorem~3.4 identifies the rational Witt target as the sheaf being
evaluated; Proposition~4.3 explains why the target section has local
preimages; Proposition~4.5, after the domain-refinement lemma, makes the
preimage on \(B\) unique; and Example~4.4 supplies the first-index
nonvanishing control.  Corollary~4.6 is not used as a premise at any stage.
It is tested only after the independent descent calculation has been
completed.

\section{Scope, controls, and conclusion}\label{sec:scope}

The index \(N=1\) is a sharp elementary control: \(V_1\) is the identity
and is lifted by \(\operatorname{id}_{\Zsh}\).  For every \(N>1\), the
prime-power part \(q^a\) of \(N\) creates a nilpotent kernel class, while
the prime-to-\(q\) part \(d\) is handled by the full set of \(d\)-th roots
of unity.  This explains why the same construction covers all composite
patterns rather than only prime indices.

There is no contradiction with the positive comparator in Deninger's
Corollary~4.7 \cite[p.~24]{Deninger2025Rational}.  That result concerns
certain finer, non-subcanonical topologies on which the relevant map
becomes an isomorphism.  Our detector and our obstruction are explicitly
formulated for the subcanonical fppf and finite-flat topologies.  A change
to a non-subcanonical topology can annihilate exactly the nilpotent
information used here.

Several boundaries are worth recording.  First, the theorem concerns
morphisms of abelian sheaves; it neither proves nor disproves the existence
of unrelated nonlinear, derived, or topology-changing constructions.
Second, no compatible Frobenius--Verschiebung law is addressed after
Verschiebung itself fails to lift.  Third, the result is absolute over the
chosen noetherian-affine site; no relative base scheme or extension to all
affine schemes is claimed.  Fourth, the bounded source search underlying
this draft supports owner subtraction but is not a global priority or
novelty theorem.

The conclusion is therefore exact and modest.  Local root factorization
does produce the expected rational Witt section, but Dedekind injectivity
forces a local representative whose double-overlap difference survives
sheafification.  That single explicit descent failure rules out an additive lift of \(V_N\)
for every \(N>1\) on each of the two sites and, for each
\(\tau\in\{\fppf,\ff\}\), gives \(V_N^*e_\tau\ne0\).

The reusable mechanism is the conditional implication from a selected
cover-local preimage, made unique, with a nonzero overlap difference to
the absence of its global preimage, followed by the formal
pushout--pullback translation for a source section mapped to that selected
target.  The example-specific verification combines Dedekind injectivity on the
root-cover source; the rational-Witt sheaf, finite-flat-cover, and
subcanonicity inputs that validate the cover and restrictions on each site;
the full big-Witt detector on the truncated nilpotent ring; and
torsion-freeness of the sheafified reduced monoid algebra.  Within this larger proof package,
four finite algebra calculations form the explicit computational core: a
root cover, a roots-of-unity product, one tensor-product overlap, and one
truncated-polynomial specialization.  This core, together with the stated
proof inputs, is the principal verification artifact of the note.

\section*{Declarations}

\paragraph{Data and materials availability.}
No empirical data were generated or analyzed.  The proof ledger, source
locator audit, claim manifest, and compilation materials used to assemble
this draft are available from the author upon reasonable request.

\paragraph{Ethics statement.}
This is a theoretical mathematics study involving no human participants,
animals, personal data, or field interventions; ethics approval and
informed consent are not applicable.

\paragraph{Author contributions.}
Liang Wang conceived the study, developed and verified the proofs, conducted the literature review, and wrote and revised the manuscript.

\paragraph{Funding.}
The author received no specific funding for this work.

\paragraph{Competing interests.}
The author declares no competing interests.

\paragraph{AI-use disclosure.}
AI-assisted tools were used during literature triage, proof-audit support,
and drafting.  Every mathematical claim, reference, attribution, and
wording choice requires final verification and responsibility by the named
human author before dissemination.  No AI system is listed as an author.

\paragraph{Limitations.}
The result is limited to additive sheaf lifts of \(V_N\) on the fixed
absolute noetherian-affine fppf and finite-flat sites.  It makes no
venue-readiness claim, no claim about all affine schemes without a newly
sheafified target, and no claim about future versions of the cited
preprint.

\bibliographystyle{plainnat}
\bibliography{references}

\end{document}
